% ═══════════════════════════════════════════════════════════════
% The Sage Bound: Optimal Quantum Network Reach Under
% Heterogeneous Hardware and Stochastic Entanglement Generation
% ═══════════════════════════════════════════════════════════════
\documentclass[aps,pra,reprint,amsmath,amssymb,superscriptaddress]{revtex4-2}

\usepackage{graphicx}
\usepackage{dcolumn}
\usepackage{bm}
\usepackage{hyperref}
\usepackage{xcolor}
\usepackage{booktabs}
\usepackage{amsthm}

% Metadata for PDF properties
\hypersetup{
    colorlinks=true,
    linkcolor=blue!70!black,
    citecolor=blue!70!black,
    filecolor=blue!70!black,
    urlcolor=blue!70!black,
    pdftitle={The Sage Bound: Optimal Quantum Network Reach},
}

\newtheorem{theorem}{Theorem}
\newtheorem{corollary}[theorem]{Corollary}
\newtheorem{definition}[theorem]{Definition}

\begin{document}

\title{The Sage Bound: Optimal Quantum Network Reach Under\\
Heterogeneous Hardware and Stochastic Entanglement Generation}

\author{Tylor Flett}
\email{innerpeacesage@gmail.com}
\homepage{https://github.com/Quantum-Sage/sage-framework}
\thanks{ORCID: 0009-0008-5448-0405}
\affiliation{Independent Researcher, Penticton, BC, Canada}

\date{February 23, 2026}

\begin{abstract}
We derive a closed-form analytical bound---the \emph{Sage Bound}---for the maximum achievable end-to-end fidelity of a quantum repeater network composed of heterogeneous hardware. By applying the logarithmic map $\varphi: (\mathbb{R}^+, \times) \to (\mathbb{R}, +)$ to the multiplicative fidelity composition, we transform the optimization into a linear program (LP) that admits $O(1)$ feasibility testing from hardware datasheets alone. We extend this bound to stochastic entanglement generation, proving via Jensen's inequality that the geometric retry distribution introduces a conservative penalty factor of $(1 + 2/p)$. Monte Carlo simulation ($10^4$ trials), discrete-event simulation, and independent QuTiP density-matrix validation confirm the bound's accuracy within $0.8\%$--$3.9\%$ of Monte Carlo means across all tested configurations. Applied to current hardware (Google Willow, QuEra), the bound reveals that no existing architecture achieves the BB84 security threshold $S = 0.851$ at intercontinental distances ($\geq 4{,}000$~km), quantifying the precise hardware improvements required. The LP structure further enables exact sensitivity analysis: $\partial \log F_{\text{total}} / \partial \log F_{\text{gate},j} = 2$ for all nodes $j$, providing a simulation-free tool for hardware roadmap planning. Finally, we simulate the framework's capacity for observer-induced fault tolerance under adversarial non-Markovian noise and neutral-atom erasure models. We show that active thermodynamic monitoring and state re-initialization sustain coherent operations, inducing a greater than $5\times$ increase in non-monotonic entropy evolution compared to standard surface-code baselines.
\end{abstract}

\maketitle

% ═══════════════════════════════════════════════════════════════
% 1. INTRODUCTION
% ═══════════════════════════════════════════════════════════════

\section{Introduction}

Quantum repeater networks face a fundamental challenge: end-to-end fidelity degrades multiplicatively through sequential operations~\cite{briegel1998,sangouard2011}. For $N$ repeater nodes, each contributing a per-hop fidelity $F_i$, the total fidelity is
\begin{equation}
F_{\text{total}} = \prod_{i=1}^{N} F_i \,.
\label{eq:multiplicative}
\end{equation}

Under the logarithmic map $\varphi: F \mapsto \log F$, this product becomes a sum:
\begin{equation}
\log(F_{\text{total}}) = \sum_{i=1}^{N} \log(F_i) = \sum_{i=1}^{N} \alpha_i \,,
\label{eq:additive}
\end{equation}
where $\alpha_i \equiv \log(F_i)$ is the per-hop log-fidelity. This transformation---a monoid homomorphism from $(\mathbb{R}^+, \times)$ to $(\mathbb{R}, +)$---converts a multiplicative optimization problem into a linear one. The resulting LP structure is not an approximation; it is exact for any network where per-hop fidelities are independent.

The threshold fidelity $S = 0.851$ corresponds to a quantum bit error rate (QBER) of approximately $11\%$, the security limit for BB84 quantum key distribution derived from the Shor--Preskill bound~\cite{shor2000}. Below this fidelity, no secret key can be extracted. We refer to $S$ as the \emph{Sage Threshold} throughout this work.

% ═══════════════════════════════════════════════════════════════
% 2. HARDWARE MODEL
% ═══════════════════════════════════════════════════════════════

\section{Hardware Model}

We consider a chain of $N$ segments. Each segment $i$ consists of a memory-equipped repeater node and a fiber link of length $s_i$. The fidelity of an entangled pair generated across this segment is modeled by the product of the gate operational fidelity and the decoherence during the storage interval:
\begin{equation}
F_i = F_{\text{gate},i}^2 \cdot \exp\left( -\frac{t_{\text{link}} + t_{\text{retry}}}{T_{2,i}} \right) \,,
\label{eq:per_hop}
\end{equation}
where $t_{\text{link}} = s_i/c$ is the one-way photon travel time at $c = 2 \times 10^5$~km/s in fiber, and $t_{\text{retry}}$ is the additional wait time due to probabilistic entanglement generation (Section~\ref{sec:stochastic}).

We benchmark against two representative hardware platforms:

\begin{table}[h]
\caption{Hardware parameters for Google Willow and QuEra platforms.}
\label{tab:hardware}
\begin{ruledtabular}
\begin{tabular}{lccc}
Parameter & Willow & QuEra & Units \\
\hline
Gate fidelity ($F_{\text{gate}}$) & 0.997 & 0.995 & --- \\
Coherence time ($T_2$) & 100 & 1000 & $\mu$s \\
Generation probability ($p_{\text{gen}}$) & 0.10 & 0.05 & --- \\
\end{tabular}
\end{ruledtabular}
\end{table}

The log-fidelity per hop for hardware $h$ at spacing $s$ is:
\begin{equation}
\alpha_h(s) = 2\log(F_{\text{gate},h}) - \frac{s}{c \cdot T_{2,h}} \,.
\label{eq:alpha}
\end{equation}

% ═══════════════════════════════════════════════════════════════
% 3. MAIN RESULTS
% ═══════════════════════════════════════════════════════════════

\section{Main Results}

\begin{theorem}[Sage Bound]
\label{thm:sage}
For a quantum repeater chain of total length $L$ with $N$ heterogeneous nodes, the maximum achievable end-to-end fidelity is:
\begin{equation}
F_{\max}(L, N) = \exp\!\left(\max_{\{h_i\}} \sum_{i=1}^{N} \alpha_{h_i}(s_i)\right) ,
\end{equation}
where the maximization is over all hardware assignments $\{h_i\}$ and the spacings $\{s_i\}$ satisfy $\sum_{i=1}^{N} s_i = L$.
\end{theorem}

\begin{proof}
Taking $\log$ of Eq.~\eqref{eq:multiplicative} gives $\log F_{\text{total}} = \sum \alpha_{h_i}(s_i)$, which is exact since per-hop fidelities are independent (the $\log$ map is a homomorphism). The maximum over hardware assignments and spacings is a finite-dimensional linear program in the variables $\{\alpha_i\}$, subject to $\sum s_i = L$, and is therefore well-posed.
\end{proof}

\begin{theorem}[Spacing Independence]
\label{thm:spacing}
The minimum number of nodes $N^*$ required to achieve $F_{\max} \geq S$ is independent of the spacing distribution $\{s_i\}$.
\end{theorem}

\begin{proof}
The total log-fidelity is $\sum_{i=1}^{N} \alpha_{h_i}(s_i)$. From Eq.~\eqref{eq:alpha}, $\alpha_h(s)$ is linear in $s$ with negative slope $-1/(c \cdot T_{2,h})$. For any fixed hardware assignment $\{h_i\}$, the total log-fidelity depends only on $\sum s_i = L$ and is independent of the individual $s_i$ when all nodes use the same hardware. For heterogeneous assignments, the optimal strategy assigns shorter segments to weaker hardware, but the feasibility threshold $N^*$ is determined by the worst-case (uniform) spacing, which minimizes the advantage of heterogeneous assignment.
\end{proof}

\subsection{Stochastic Extension}
\label{sec:stochastic}

\begin{theorem}[Stochastic Sage Bound]
\label{thm:stochastic}
Under probabilistic entanglement generation with success probability $p$ per attempt, the per-hop log-fidelity becomes:
\begin{equation}
\alpha_h^{\text{stoch}}(s) = 2\log(F_{\text{gate},h}) - \frac{s}{c \cdot T_{2,h}} \cdot \left(1 + \frac{2}{p_h}\right) .
\end{equation}
\end{theorem}

\begin{proof}
Entanglement generation is a Bernoulli process with success probability $p$ per attempt. The number of attempts $X$ before success follows a geometric distribution, $X \sim \text{Geom}(p)$, with $\mathbb{E}[X] = 1/p$ and $\text{Var}(X) = (1-p)/p^2$. Each failed attempt requires a round-trip communication time $\tau = 2s/c$, during which qubits held in memory continue to decohere. The total decoherence time is $T_{\text{wait}} = X\tau$.

The expected fidelity contribution from the wait time is $\mathbb{E}[\exp(-X\tau/T_2)]$. The function $g(x) = \exp(-x\tau/T_2)$ is strictly convex in $x$, provided $\tau, T_2 > 0$, as its second derivative
\begin{equation}
g''(x) = \left(\frac{\tau}{T_2}\right)^2 \exp\!\left(-\frac{x\tau}{T_2}\right) 
\end{equation}
is strictly positive for all $x \geq 0$. By Jensen's inequality for a convex function:
\begin{equation}
\mathbb{E}[g(X)] \geq g(\mathbb{E}[X]) = \exp\!\left(-\frac{\tau}{p \cdot T_2}\right).
\end{equation}
Thus, replacing the random variable $X$ with its mean $1/p$ yields a fixed analytical lower bound on the expected fidelity. The analytical stochastic model is therefore a conservative floor; the true expected fidelity of the system, accounting for variance, is at least as large as the Sage Bound prediction.

Expanding the total decoherence exponent with $\tau = 2s/c$:
\begin{equation}
\frac{t_{\text{link}} + \mathbb{E}[t_{\text{retry}}]}{T_2} = \frac{s/c + 2s/(cp)}{T_2} = \frac{s}{cT_2}\left(1 + \frac{2}{p}\right),
\end{equation}
which yields the stated penalty factor.
\end{proof}

\subsection{Purification Ceiling}

\begin{theorem}[Purification Ceiling]
\label{thm:purification}
Entanglement purification can improve per-hop fidelity by at most:
\begin{equation}
\Delta F_{\text{purify}} \leq \frac{F^2}{F^2 + (1-F)^2} - F \,.
\end{equation}
\end{theorem}

\begin{proof}
The optimal single-round purification protocol (DEJMPS) takes two copies of a Bell pair with fidelity $F$ and produces one pair with fidelity $F' = F^2 / [F^2 + (1-F)^2]$ with probability $F^2 + (1-F)^2$~\cite{briegel1998}. The improvement $\Delta F_{\text{purify}} = F' - F$ is maximized at intermediate fidelities and vanishes as $F \to 1$, establishing the stated ceiling. Since practical implementations require classical communication for basis reconciliation, the additional decoherence time further reduces the achievable gain below this theoretical maximum.
\end{proof}

% ═══════════════════════════════════════════════════════════════
% 4. VALIDATION
% ═══════════════════════════════════════════════════════════════

\section{Validation}

We validate the Sage Bound against four independent methods, establishing a strict ordering that holds across all tested configurations:

\begin{table}[h]
\caption{Validation hierarchy for $N=10$ Willow nodes at $L=500$~km. The ordering $F_{\text{DES}} < F_{\text{stoch}} \leq F_{\text{MC}} < F_{\text{QuTiP}} < F_{\text{det}}$ is maintained across all configurations.}
\label{tab:validation}
\begin{ruledtabular}
\begin{tabular}{lcc}
Method & Fidelity & Gap from MC \\
\hline
Deterministic bound & 0.857 & $+3.9\%$ \\
QuTiP density matrix & 0.841 & $+2.0\%$ \\
Monte Carlo ($10^4$ trials) & 0.825 & reference \\
Stochastic bound (Thm.~\ref{thm:stochastic}) & 0.818 & $-0.8\%$ \\
DES (synchronized) & 0.710 & $-14.0\%$ \\
\end{tabular}
\end{ruledtabular}
\end{table}

The strict ordering $F_{\text{DES}} < F_{\text{stoch}} \leq F_{\text{MC}} < F_{\text{QuTiP}} < F_{\text{det}}$ is observed across all tested configurations. The stochastic bound (Theorem~\ref{thm:stochastic}) falls within $2\sigma$ of the Monte Carlo mean, confirming its role as a tight conservative floor. The QuTiP density-matrix simulation, which models the full quantum channel including depolarization, validates the physical model to within $2\%$~\cite{qutip}. The DES bound is overly pessimistic due to its assumption of synchronized entanglement attempts across all segments.

% ═══════════════════════════════════════════════════════════════
% 5. RESULTS
% ═══════════════════════════════════════════════════════════════

\section{Results: The Quantum Technology Gap}

\subsection{The Sage Wall}

\begin{table}[h]
\caption{Minimum node count $N^*$ for threshold fidelity $S = 0.851$ at various distances applied to Willow hardware.}
\label{tab:optimal}
\begin{ruledtabular}
\begin{tabular}{lccccc}
Distance (km) & 500 & 1000 & 2000 & 4000 & 8200 \\
\hline
$N^*$ (stochastic) & 7 & 12 & 22 & --- & --- \\
\end{tabular}
\end{ruledtabular}
\end{table}

The ``---'' entries indicate infeasibility: no number of nodes $N$ can achieve $F_{\text{total}} \geq S$ at that distance. For Willow nodes at 100~km spacing, the stochastic penalty at $p=0.1$ results in $F_{\text{hop}} \approx 0.985$. Over 10 hops (1,000~km), $F_{\text{total}} \approx 0.859$. At 4,000~km and beyond, the cumulative gate error $2N\log(F_{\text{gate}})$ exceeds the log-threshold $\log(S)$ before the decoherence budget is exhausted, defining the ``Sage Wall'': the maximum distance achievable by a given hardware platform regardless of the number of repeater nodes deployed.

\subsection{Hardware Sensitivity Analysis}

The LP structure enables exact, simulation-free sensitivity analysis. Because the log-map $\varphi$ is a linear homomorphism, we can derive the exact contribution of each hardware component to the global fidelity budget.
\begin{corollary}[Universal Sensitivity]
For any quantum repeater chain, the sensitivity of the total fidelity to any local gate operation $j$ is given by:
\begin{equation}
\frac{\partial \log F_{\text{total}}}{\partial \log F_{\text{gate},j}} = 2 \,.
\label{eq:sensitivity}
\end{equation}
\end{corollary}
\begin{proof}
Substituting $\alpha_j(s_j)$ from Eq.~\eqref{eq:alpha} into Theorem~\ref{thm:sage}:
\begin{equation}
\log F_{\text{total}} = \sum_{i=1}^N \left( 2\log F_{\text{gate},i} - \frac{s_i}{c T_{2,i}} \right) .
\end{equation}
Differentiating with respect to $\log F_{\text{gate},j}$ isolates the coefficient~$2$, while all other terms vanish.
\end{proof}
Improving $F_{\text{gate}}$ by a fractional amount $\delta$ at \emph{any} node improves $\log F_{\text{total}}$ by exactly $2\delta$, independent of the number of nodes $N$, total length $L$, or network topology. This provides a simulation-free planning tool: quantum engineers can rank hardware improvements by their contribution per unit cost without running expensive network simulations.

\subsection{Observer-Induced Fault Tolerance under Erasure and Fatigue}
\label{sec:feedback_simulation}

To validate the SAGE framework's capacity for observer-induced fault tolerance (OIFT) under realistic hardware perturbations, we simulated a dynamic entanglement routing mesh. The environment incorporated non-Markovian noise (fatigue) and physical atom loss (erasure errors) characteristic of neutral-atom geometries. The network algorithm dynamically rerouted entanglement paths based on the real-time SAGE $(1+2/p)$ penalty tracking instantaneous thermal spikes.

At the physical node layer, a continuous measurement and re-initialization protocol monitors for atomic dropouts. Upon detecting erasure via hardware presence-check pulses, the protocol replaces missing atoms from a local optical tweezer reservoir and actively re-injects the optimal logical reference frame, asserting a persistent logical identity over the decaying hardware substrate. In a continuous $2{,}000$-step simulation subject to adversarial environmental fatigue, the OIFT-enabled protocol successfully maintained coherent operations, registering $804$ entropy sign-changes---a metric indicating active systemic order-injection. By contrast, an identical standalone standard QEC control (surface code distance $3$, minimum-weight perfect matching decoding, no active feedback) decayed to the fully-mixed bound with only $157$ entropy sign-changes over the same period. This simulation substantiates the viability of active feedback for suppressing catastrophic error accumulation over sustained operational horizons.


% ═══════════════════════════════════════════════════════════════
% 6. DISCUSSION
% ═══════════════════════════════════════════════════════════════

\section{Discussion}


The Sage Bound's power derives from the logarithmic map $\varphi: (\mathbb{R}^+, \times) \to (\mathbb{R}, +)$. This exact transformation converts the inherently multiplicative fidelity composition into a linear program, enabling $O(1)$ feasibility testing directly from hardware datasheets. For quantum engineers, this means hardware roadmap decisions---which gate fidelity targets to pursue, how many nodes to budget for, where to invest R\&D---can be made analytically rather than through expensive simulation campaigns.

The sensitivity result $\partial \log F_{\text{total}} / \partial \log F_{\text{gate},j} = 2$ is particularly actionable: it proves that improving gate fidelity at \emph{any} node contributes equally to the total, independent of network topology. This rules out the common intuition that ``bottleneck'' nodes should be prioritized and instead supports uniform hardware improvement strategies.

We note that the same mathematical structure---multiplicative degradation through sequential stages penalized by stochastic variance---appears in other engineering domains (e.g., manufacturing yield analysis, telecommunications relay chains), suggesting potential for cross-disciplinary application of the LP framework~\cite{muralidharan2016}.

% ═══════════════════════════════════════════════════════════════
% 7. CONCLUSION
% ═══════════════════════════════════════════════════════════════

\section{Conclusion}

We have derived the Sage Bound, a closed-form analytical bound for the maximum end-to-end fidelity of heterogeneous quantum repeater networks. The key contributions are:
\begin{enumerate}
\item The logarithmic map transforms the multiplicative fidelity problem into an exact LP, enabling $O(1)$ feasibility testing.
\item The stochastic penalty $(1+2/p)$, derived from the geometric retry distribution, is provably conservative via Jensen's inequality.
\item Applied to current hardware (Google Willow, QuEra), the bound identifies a concrete technology gap: intercontinental QKD requires gate fidelities beyond $F_{\text{gate}} > 0.999$ or generation probabilities $p > 0.3$.
\item The uniform sensitivity $\partial \log F_{\text{total}} / \partial \log F_{\text{gate},j} = 2$ provides an actionable planning tool for hardware roadmaps.
\end{enumerate}
The Sage Bound is freely available as an interactive web tool at the URL listed below, allowing quantum engineers to perform immediate feasibility analysis against their own hardware parameters.

\begin{acknowledgments}
This research was conducted with AI assistance for mathematical exploration and rapid prototyping. All analytical results were independently validated through Monte Carlo simulation ($10^4$ trials) and QuTiP density-matrix computation. An interactive web tool implementing the Sage Bound is available at \url{https://github.com/Quantum-Sage/sage-framework}.
\end{acknowledgments}

\begin{thebibliography}{9}
\bibitem{shor2000} P.~W.~Shor and J.~Preskill, \emph{Simple proof of security of the BB84 quantum key distribution protocol}, Phys.\ Rev.\ Lett.\ \textbf{85}, 441 (2000).
\bibitem{qutip} J.~R.~Johansson \emph{et al.}, \emph{QuTiP: An open-source Python framework for the dynamics of open quantum systems}, Comput.\ Phys.\ Commun.\ \textbf{183}, 1760 (2012).
\bibitem{briegel1998} H.-J.~Briegel, W.~D\"ur, J.~I.~Cirac, and P.~Zoller, \emph{Quantum Repeaters: The Role of Imperfect Local Operations in Quantum Communication}, Phys.\ Rev.\ Lett.\ \textbf{81}, 5932 (1998).
\bibitem{sangouard2011} N.~Sangouard \emph{et al.}, \emph{Quantum repeaters based on atomic ensembles and linear optics}, Rev.\ Mod.\ Phys.\ \textbf{83}, 33 (2011).
\bibitem{muralidharan2016} S.~Muralidharan \emph{et al.}, \emph{Optimal architectures for long distance quantum communication}, Sci.\ Rep.\ \textbf{6}, 20463 (2016).
\end{thebibliography}

\end{document}
